
\documentclass[a4paper,10pt]{article}

\usepackage[utf8]{inputenc}
\usepackage[T1]{fontenc}
\usepackage[turkish]{babel}
\usepackage{geometry}
\usepackage{booktabs}
\usepackage{array}
\usepackage{tabularx}
\usepackage{longtable}
\usepackage{enumitem}
\usepackage{listings}
\usepackage{xcolor}
\usepackage[hidelinks]{hyperref}
\usepackage{microtype}
\usepackage{amsmath}
\usepackage{titlesec}

% Türkçe babel paketinin "=" karakterini listings seçeneklerinde etkilemesini engeller
\AtBeginDocument{\shorthandoff{=}}

\geometry{
    a4paper,
    top=1.8cm,
    bottom=1.8cm,
    left=1.8cm,
    right=1.8cm
}

\setlength{\parindent}{0pt}
\setlength{\parskip}{5pt}
\renewcommand{\arraystretch}{1.18}

\titleformat{\section}{\large\bfseries}{\thesection.}{0.5em}{}
\titleformat{\subsection}{\normalsize\bfseries}{\thesubsection.}{0.5em}{}

\lstdefinestyle{cstyle}{
    language={[ANSI]C},
    basicstyle=\ttfamily\small,
    keywordstyle=\bfseries,
    commentstyle=\itshape,
    numbers=left,
    numberstyle=\tiny,
    numbersep=7pt,
    frame=single,
    breaklines=true,
    showstringspaces=false,
    tabsize=4
}

\newcommand{\hex}[1]{\texttt{0x#1}}
\newcommand{\code}[1]{\texttt{#1}}

\title{\textbf{DEOS Data Types \& Config}\\
\large Sürüm 0.3}
\author{DEOS -- Dokuz Eylül Otonom Sistemler}
\date{}

\begin{document}
\maketitle

\begin{center}
\begin{tabular}{ll}
\textbf{Doküman tipi:} & Interface Control Document (ICD) \\
\textbf{Protokol sürümü:} & v1.1 (\hex{11}) \\
\textbf{Taşıma altyapısı:} & CAN-FD ana omurga \\
\textbf{Durum:} & ICD Sürüm 0.3 -- v1.1 wire-format revizyonu \\
\end{tabular}
\end{center}


\section{Veri Tipleri ve Byte Sırası}

DEOS protokolünde çok baytlı alanlar için varsayılan byte sırası \textbf{Little-Endian}
olarak tanımlanır.

Wire formatında temel tercih, doğrudan IEEE-754 floating-point taşımak yerine ölçekli
integer değer kullanmaktır. Böylece bir sinyalin fiziksel karşılığı açık şekilde
\code{raw value + scale + offset + unit} ile tanımlanır.

Genel fiziksel değer dönüşümü:

\[
PhysicalValue = RawValue \times Scale + Offset
\]

\subsection{Normatif Veri Tipleri}

\begin{center}
\begin{tabularx}{\textwidth}{l c X}
\toprule
Tip & Boyut & Kullanım \\
\midrule
\code{uint8\_t}  & 1 byte & Enum, yüzde, küçük unsigned değerler \\
\code{int8\_t}   & 1 byte & Küçük signed değerler \\
\code{uint16\_t} & 2 byte & Ölçekli pozitif fiziksel değerler \\
\code{int16\_t}  & 2 byte & Ölçekli signed fiziksel değerler \\
\code{uint32\_t} & 4 byte & Sayaç / geniş pozitif değerler \\
\code{int32\_t}  & 4 byte & Geniş signed fiziksel değerler \\
\code{float32}   & 4 byte & Yalnız ilgili mesaj tanımında açıkça belirtilirse IEEE-754 \\
\bottomrule
\end{tabularx}
\end{center}

Floating-point kullanımına genel yasak konulmamıştır; ancak bir alan \code{float32}
olarak tanımlanmamışsa integer tabanlı ölçekleme varsayılır.

\subsection{Başlangıç Fiziksel Sinyal Formatları}

\begin{center}
\begin{tabularx}{\textwidth}{l l l l X}
\toprule
Sinyal & Wire Type & Scale & Unit & Örnek \\
\midrule
Steering target angle & \code{int16\_t} & 0.01 & degree & 1250 = 12.50 deg \\
Vehicle target speed & \code{uint16\_t} & 0.01 & m/s & 725 = 7.25 m/s \\
Brake target & \code{uint16\_t} & 0.01 & \% & 7500 = 75.00\% \\
BMS voltage & \code{uint16\_t} & 0.01 & V & 7425 = 74.25 V \\
BMS current & \code{int16\_t} & 0.1 & A & -125 = -12.5 A \\
BMS SOC & \code{uint16\_t} & 0.01 & \% & 8530 = 85.30\% \\
Temperature & \code{int16\_t} & 0.1 & degC & 253 = 25.3 degC \\
\bottomrule
\end{tabularx}
\end{center}

Bu tablo başlangıç wire-format standardını tanımlar. Nihai min/max sınırlar ve
invalid değerler ilgili komut veya sinyal tanımında ayrıca belirtilmelidir.

\subsection{Ortak Parameter Command ID'leri}

Aşağıdaki command ID'leri configurable servislerde ortak olarak kullanılır:

\begin{center}
\begin{tabular}{c l}
\toprule
ID & Symbol \\
\midrule
\hex{E0} & \code{DEOS\_CMD\_GET\_PARAMETER} \\
\hex{E1} & \code{DEOS\_CMD\_SET\_PARAMETER} \\
\hex{E2} & \code{DEOS\_CMD\_SAVE\_PARAMETERS} \\
\hex{E3} & \code{DEOS\_CMD\_RESTORE\_DEFAULTS} \\
\hex{E4} & \code{DEOS\_CMD\_GET\_PARAMETER\_INFO} \\
\bottomrule
\end{tabular}
\end{center}

\begin{lstlisting}[style=cstyle]
typedef enum
{
    DEOS_CMD_GET_PARAMETER       = 0xE0,
    DEOS_CMD_SET_PARAMETER       = 0xE1,
    DEOS_CMD_SAVE_PARAMETERS     = 0xE2,
    DEOS_CMD_RESTORE_DEFAULTS    = 0xE3,
    DEOS_CMD_GET_PARAMETER_INFO  = 0xE4
} deos_parameter_command_t;
\end{lstlisting}

\subsection{Parameter ID Alanı}

Parameter ID 16-bit olarak tanımlanır. Parameter ID global değildir;
\code{Service ID + Parameter ID} birlikte benzersizdir.

Örnek Steering Parameter Registry:

\begin{center}
\begin{tabularx}{\textwidth}{c l X}
\toprule
ID & Symbol & Açıklama \\
\midrule
\hex{0001} & \code{DEOS\_PARAM\_STEERING\_PID\_KP} & PID oransal katsayısı \\
\hex{0002} & \code{DEOS\_PARAM\_STEERING\_PID\_KI} & PID integral katsayısı \\
\hex{0003} & \code{DEOS\_PARAM\_STEERING\_PID\_KD} & PID türev katsayısı \\
\hex{0010} & \code{DEOS\_PARAM\_STEERING\_ZERO\_OFFSET} & Direksiyon sıfır ofseti \\
\hex{0011} & \code{DEOS\_PARAM\_STEERING\_MIN\_ANGLE} & Minimum izinli açı \\
\hex{0012} & \code{DEOS\_PARAM\_STEERING\_MAX\_ANGLE} & Maksimum izinli açı \\
\hex{0013} & \code{DEOS\_PARAM\_STEERING\_CURRENT\_LIMIT} & Akım limiti \\
\bottomrule
\end{tabularx}
\end{center}

Örnek Traction Parameter Registry:

\begin{center}
\begin{tabularx}{\textwidth}{c l X}
\toprule
ID & Symbol & Açıklama \\
\midrule
\hex{0001} & \code{DEOS\_PARAM\_TRACTION\_PID\_KP} & PID oransal katsayısı \\
\hex{0002} & \code{DEOS\_PARAM\_TRACTION\_PID\_KI} & PID integral katsayısı \\
\hex{0003} & \code{DEOS\_PARAM\_TRACTION\_PID\_KD} & PID türev katsayısı \\
\hex{0010} & \code{DEOS\_PARAM\_TRACTION\_MAX\_CURRENT} & Maksimum akım \\
\hex{0011} & \code{DEOS\_PARAM\_TRACTION\_MAX\_SPEED} & Maksimum hız \\
\hex{0012} & \code{DEOS\_PARAM\_TRACTION\_ACCEL\_LIMIT} & İvme limiti \\
\bottomrule
\end{tabularx}
\end{center}

Örnek Brake Parameter Registry:

\begin{center}
\begin{tabularx}{\textwidth}{c l X}
\toprule
ID & Symbol & Açıklama \\
\midrule
\hex{0001} & \code{DEOS\_PARAM\_BRAKE\_PID\_KP} & Gelecek kullanım / gerekli olması halinde \\
\hex{0002} & \code{DEOS\_PARAM\_BRAKE\_PID\_KI} & Gelecek kullanım / gerekli olması halinde \\
\hex{0003} & \code{DEOS\_PARAM\_BRAKE\_PID\_KD} & Gelecek kullanım / gerekli olması halinde \\
\hex{0010} & \code{DEOS\_PARAM\_BRAKE\_MIN\_POSITION} & Minimum aktüatör konumu \\
\hex{0011} & \code{DEOS\_PARAM\_BRAKE\_MAX\_POSITION} & Maksimum aktüatör konumu \\
\hex{0012} & \code{DEOS\_PARAM\_BRAKE\_CURRENT\_LIMIT} & Akım limiti \\
\bottomrule
\end{tabularx}
\end{center}

\section{State ve Mode Registry}

State, düğümün anlık çalışma durumunu; Mode ise düğümün hangi çalışma davranışı altında
olduğunu ifade eder. İki alan birbirinden bağımsızdır.

\subsection{State Enum}

\begin{center}
\begin{tabular}{c l}
\toprule
Değer & Symbol \\
\midrule
\hex{00} & \code{DEOS\_STATE\_INIT} \\
\hex{01} & \code{DEOS\_STATE\_STANDBY} \\
\hex{02} & \code{DEOS\_STATE\_READY} \\
\hex{03} & \code{DEOS\_STATE\_ACTIVE} \\
\hex{04} & \code{DEOS\_STATE\_CALIBRATING} \\
\hex{05} & \code{DEOS\_STATE\_SAFE} \\
\hex{06} & \code{DEOS\_STATE\_FAULT} \\
\hex{FF} & \code{DEOS\_STATE\_UNKNOWN} \\
\bottomrule
\end{tabular}
\end{center}

\begin{lstlisting}[style=cstyle]
typedef enum
{
    DEOS_STATE_INIT        = 0x00,
    DEOS_STATE_STANDBY     = 0x01,
    DEOS_STATE_READY       = 0x02,
    DEOS_STATE_ACTIVE      = 0x03,
    DEOS_STATE_CALIBRATING = 0x04,
    DEOS_STATE_SAFE        = 0x05,
    DEOS_STATE_FAULT       = 0x06,
    DEOS_STATE_UNKNOWN     = 0xFF
} deos_state_t;
\end{lstlisting}

\subsection{Mode Enum}

\begin{center}
\begin{tabular}{c l}
\toprule
Değer & Symbol \\
\midrule
\hex{00} & \code{DEOS\_MODE\_NORMAL} \\
\hex{01} & \code{DEOS\_MODE\_MANUAL} \\
\hex{02} & \code{DEOS\_MODE\_AUTONOMOUS} \\
\hex{03} & \code{DEOS\_MODE\_SERVICE} \\
\hex{04} & \code{DEOS\_MODE\_TEST} \\
\hex{FF} & \code{DEOS\_MODE\_UNKNOWN} \\
\bottomrule
\end{tabular}
\end{center}

\begin{lstlisting}[style=cstyle]
typedef enum
{
    DEOS_MODE_NORMAL     = 0x00,
    DEOS_MODE_MANUAL     = 0x01,
    DEOS_MODE_AUTONOMOUS = 0x02,
    DEOS_MODE_SERVICE    = 0x03,
    DEOS_MODE_TEST       = 0x04,
    DEOS_MODE_UNKNOWN    = 0xFF
} deos_mode_t;
\end{lstlisting}

\end{document}
